\documentclass[10pt,a4paper]{amsart}
\usepackage[margin=19mm]{geometry}
\usepackage{amsmath,amssymb,mathtools}
\usepackage[scr=rsfs,cal=euler]{mathalfa}
\usepackage{xfrac}
\usepackage[unicode,hidelinks,bookmarks=false]{hyperref}
\newcommand{\ie}{i.e.\ }
\newcommand{\cf}{cf.\ }
\newcommand{\resp}{respectively\ }
\newcommand{\Subsection}{\subsection}
\providecommand{\nobreakdash}{\nobreak\hbox{-}}
\setlength{\emergencystretch}{2em}
\multlinegap0pt
\usepackage{xcolor}
\usepackage{amssymb}
\usepackage{mathtools}
\usepackage{dsfont}
\usepackage{enumerate}
\usepackage[shortlabels]{enumitem}
\usepackage{tikz}
\usepackage{braket}
\newtheorem{theorem}{Theorem}
\newtheorem{corollary}{Corollary}
\newcommand\bC{{\mathbb C}}
\newcommand\cG{{\mathcal G}}
\newcommand\cO{{\mathcal O}}
\newcommand{\tr}{\text{tr}} % trace v2
\title{Local quantum Cuntz--Krieger algebras of dephasing quantum graphs}
\author{Lara Ismert}
\date{Source: arXiv:2608.23520v1 (CC BY 4.0)}
\begin{document}
\maketitle

\noindent\emph{Source and licence.} Local quantum Cuntz--Krieger algebras of dephasing quantum graphs. Version arXiv:2608.23520v1 (CC BY 4.0), distributed by arXiv under the Creative Commons Attribution 4.0 International licence, http://creativecommons.org/licenses/by/4.0/, as recorded in the arXiv licence metadata for this exact version and shown on the abstract page https://arxiv.org/abs/2608.23520v1. Full licence notice: \texttt{licenses/arxiv-2608.23520-CC-BY-4.0.txt}.\par\smallskip

\noindent\textit{Editorial scope.} Counted unit: the article's corollary cor:cuntz\_isometries\_for\_L(M\_n-tr-D (source0.tex lines 968--975). The article's complete original proof of it is delivered with it (source0.tex lines 977--1022, one \texttt{proof} environment of 46 lines). No mathematics is written, repaired or re-derived: the statement and the proof are the article's own lines, in the article's own notation, and no retained line is substituted or re-flowed.  Retained uncounted as the source-local context the proof needs: lines 629--642; lines 905--920.  Nothing was omitted from any retained span, so the package carries no omission record.  No retained line contains a citation, so the wrapper carries no bibliography and no bibliography entry.  No retained line contains a figure environment, an image-inclusion command or drawing code, so no image is delivered and the package declares no asset dependency.  Editorial formatting only, in addition: an AMS \texttt{amsart} preamble that restates the article's own declarations and macros used by the retained lines, namely the environments \texttt{theorem}, \texttt{corollary} on separate counters, together with the standard packages those lines need.  The retained environments print in this document as Theorem 1, Corollary 1 in order of appearance, whereas the article prints the same items with the numbering of its own full document.  The complete native source of the article as distributed in the arXiv e-print for this exact version is bundled unchanged as \texttt{sources/208414/source0.tex}.
\begin{theorem}\label{thm:lqck_structure_theorem}
    Suppose $B\cong \oplus_{a=1}^d M_{n_a}(\bC)$, and let $\cG:=(B,\psi,A)$ be a quantum graph such that $\psi$ is a $\delta$-form and $A$ is cp. Given any local quantum Cuntz--Krieger $\cG$-family $S:B\to \mathcal{C}$, for all $1\leq a,b\leq d$ and $1\leq i,j\leq n_a$, $1\leq k,\ell\leq n_b$: 
    \begin{enumerate}
        \item $(S^{(a)}_{ij})^* S^{(b)}_{k\ell} \neq 0 
        \quad \Longrightarrow \quad 
        a=b\text{ and }i=k$.
        \item $S^{(a)}_{ij} S^{(b)}_{k\ell} \neq 0 
        \quad \quad \; \Longrightarrow \quad
        f_{ij}^{(a)}\cdot \epsilon_\cG\cdot f_{k\ell}^{(b)}\neq 0$.
        \item $S^{(a)}_{ij} (S^{(b)}_{k\ell})^* \neq 0 
        \quad \,\Longrightarrow \quad
        \exists 1\leq m\leq n_b$ such that $f_{ij}^{(a)}\cdot \epsilon_\cG\cdot A(f_{\ell m}^{(b)}) \neq 0$. 
    \end{enumerate}
\end{theorem}

\begin{theorem}\label{thm:unital_embedding_of_O_n}
     Suppose $B\cong \oplus_{a=1}^d M_n$ for some fixed $n\in \mathbb{N}$, and let $s:B\to L\cO(B,\psi,D)$ denote the universal local quantum Cuntz--Krieger $(B,\psi,D)$-family. For each $1\leq a\leq d$ and $1\leq i,j \leq n_a$, set 
     \[
     t^{(a)}_{ij}
     :=
     \delta\psi(e_{ii}^{(a)})^{1/2}\cdot s(f_{ij}^{(a)}).
     \] 
     and 
      \[
     \forall 1\leq i\leq n: \quad
     V_i
     :=\sum_{a=1}^d\sum_{j=1}^n t^{(a)}_{ij}.
     \] 
     Then the collection $\{V_i:1\leq i\leq n\}$ 
    consists of Cuntz isometries which generate a unital embedding of $\cO_n$ inside $L\cO(B,\psi,D).$
\end{theorem}

\begin{corollary}\label{cor:cuntz_isometries_for_L(M_n-tr-D}
    Let $s:M_n\to L\cO(M_n,\tr,D)$ be the universal local quantum Cuntz--Krieger $(M_n,\tr,D)$-family. For each $1\leq i,j \leq n$, set \[t_{ij}:=\sqrt{n}\cdot s(f_{ij}).\] Then the collection $\{V_i:1\leq i\leq n\}$ where 
     \[
     \forall 1\leq i\leq n: \quad
     V_i
     :=\sum_{j=1}^n t_{ij}
     \] consists of Cuntz isometries which generate all of $L\cO(M_n,\tr,D).$
\end{corollary}

\begin{proof}
Note that $\sqrt{n} = \delta \psi(e_{ii})^{1/2}$ when $\psi=\tr$. Hence, Theorem~\ref{thm:unital_embedding_of_O_n} confirms $\{V_i:1\leq i \leq n\}$ are Cuntz isometries, so it remains to show that $C^*(\{V_i:1\leq i \leq n\})=L\cO(M_n,\tr,D).$ For any $1\leq i,j,k,\ell\leq n$, as $t_{ij}t_{k\ell}\neq 0\iff 
S_{ij}S_{k\ell} \neq 0$, Theorem~\ref{thm:lqck_structure_theorem} yields
\[
t_{ij}t_{k\ell}\neq 0 
\quad
\Longrightarrow 
\quad
f_{ij}\cdot \epsilon_D \cdot f_{k\ell}\neq 0 
\quad
{\Longrightarrow}
\quad
j=k,
\]
which tells us $t_{ik}V_j = \sum_{\ell=1}^d t_{ik}t_{j\ell} = 0$ unless $j=k.$ Hence,
\[V_iV_j
%=\left(\sum_k t_{ik}\right)\left(\sum_{\ell} t_{j\ell}\right)
=\sum_{k,\ell=1}^n t_{ik}t_{j\ell}
%=\delta_{j=k}\sum_{k,\ell=1}^n t_{ik}t_{k\ell}
=\sum_{\ell=1}^n t_{ij}t_{j\ell}
=t_{ij}\left(\sum_{\ell=1}^n t_{j\ell}\right)
=t_{ij}V_j.
\]
Using 
$
{\bf 1}
=
\sum_{k=1}^d V_kV_k^*
$
from Theorem~\ref{thm:unital_embedding_of_O_n}, we have
\[
    t_{ij}
    =t_{ij}{\bf 1}
    =t_{ij}\left(\sum_{k=1}^nV_kV_k^*\right)
    =\sum_{k=1}^n t_{ij}V_kV_k^*
    =t_{ij}V_jV_j^*
    =V_iV_jV_j^*.
\] 
Therefore, all of the $t_{ij}$ belong to $C^*(\{V_i:1\leq i\leq n\})$, and because the $t_{ij}$ generate $L\cO(M_n,\tr,D)$, we obtain
$L\cO(M_n,\tr,D)
%=
%C^*(\{t_{ij}:1\leq i,j\leq n\})
%=
%C^*(\{V_i:1\leq i\leq n\})
\cong \cO_n$.%
\end{proof}

\end{document}
